\begin{helper}
Let \(C^0_{\omega_0}\) and \(C^{\prime0}_{\omega_0}\) be paired source and
target supported base cells, indexed by \(\omega_0\in\Omega_0\).  Let
\(\Lambda\) be a finite set of secondary Boolean labels, with source tests
\(Q_\lambda\) and target tests \(Q'_\lambda\).  For a truth vector
\(\tau\subseteq\Lambda\), write
\[
C^0_{\omega_0,\tau}
=C^0_{\omega_0}\cap
  \{\eta:\eta\in Q_\lambda\Longleftrightarrow \lambda\in\tau
      \text{ for every }\lambda\in\Lambda\}
\]
and define \(C^{\prime0}_{\omega_0,\tau}\) analogously using
\(C^{\prime0}_{\omega_0}\) and \(Q'_\lambda\).

Assume the concrete source atoms \(C^0_{\omega_0,\tau}\) are measurable,
cover \(C^0_{\omega_0}\), and are pairwise disjoint as \(\tau\) varies.
Assume likewise that the concrete target atoms
\(C^{\prime0}_{\omega_0,\tau}\) are measurable, cover
\(C^{\prime0}_{\omega_0}\), and are pairwise disjoint as \(\tau\) varies.
Assume also that each transported source atom and target atom have equal
measure,
\[
\nu(C^0_{\omega_0,\tau})=\nu'(C^{\prime0}_{\omega_0,\tau}),
\]
and that this common atom mass is represented by a nonnegative real weight:
there is \(w_{\omega_0,\tau}\ge0\) with
\[
\operatorname{ofReal}(w_{\omega_0,\tau})
=\nu(C^0_{\omega_0,\tau})
=\nu'(C^{\prime0}_{\omega_0,\tau}).
\]
Then the four mass hypotheses needed by the finite stage-cell common
refinement hold for these concrete cells and labels.  Namely,
\[
\nu(C^0_{\omega_0})
=\sum_{\tau\subseteq\Lambda}\nu(C^0_{\omega_0,\tau}),
\qquad
\nu'(C^{\prime0}_{\omega_0})
=\sum_{\tau\subseteq\Lambda}\nu'(C^{\prime0}_{\omega_0,\tau}),
\]
corresponding source and target atom masses are equal, and those common atom
masses have nonnegative real representatives.
\end{helper}
\begin{proof}
For each fixed \(\omega_0\), the source splitting hypothesis is the finite
additivity identity for the actual source measurable partition
\(\{C^0_{\omega_0,\tau}:\tau\subseteq\Lambda\}\).  This partition consists
of all simultaneous truth-vector atoms of the finitely many source Boolean
tests inside the supported base cell.  The atoms cover \(C^0_{\omega_0}\)
and are pairwise disjoint by hypothesis; the cover is also forced by the
truth-vector construction, since each state has the unique truth vector
consisting of the labels whose source tests contain it.  Because
\(\Lambda\) is finite, there are only finitely many truth vectors.  Finite
additivity for measures over a finite pairwise-disjoint measurable union
therefore gives
\[
\nu(C^0_{\omega_0})
=\sum_{\tau\subseteq\Lambda}\nu(C^0_{\omega_0,\tau}).
\]
This is the first output hypothesis.

The target argument is identical.  For each \(\omega_0\), the target atoms
\(\{C^{\prime0}_{\omega_0,\tau}:\tau\subseteq\Lambda\}\) are the simultaneous
truth-vector atoms of the transported target Boolean tests inside
\(C^{\prime0}_{\omega_0}\).  They cover the target base cell, are pairwise
disjoint, and are measurable by hypothesis.  Finite additivity over this
finite disjoint measurable cover gives the displayed sum for
\(\nu'(C^{\prime0}_{\omega_0})\).  This is the second output hypothesis.

The third output hypothesis is the assumed transported atom-mass equality:
for every \(\omega_0\) and truth vector \(\tau\), the source atom
\(C^0_{\omega_0,\tau}\) and target atom
\(C^{\prime0}_{\omega_0,\tau}\) have equal measure.  This equality is not
derived from the abstract common-refinement lemma or from the unrefined base
cell weight; it must come from the concrete product and extension laws for
the specific supported atoms under consideration.

The fourth output hypothesis is the assumed finite real representation of
each common atom mass.  The chosen representative
\(w_{\omega_0,\tau}\) is nonnegative and satisfies
\[
\operatorname{ofReal}(w_{\omega_0,\tau})
=\nu(C^0_{\omega_0,\tau})
=\nu'(C^{\prime0}_{\omega_0,\tau}).
\]
Thus the source splitting, target splitting, atom equality, and nonnegative
atom-weight clauses are all available as a single registered supplier for
the stage-cell common-refinement node.
\end{proof}
